%%
%% part4-advanced/ch26-large-projects.tex
%%
%% Chapter 26: Large Projects — How to organize books with
%% multiple files, parts, and contributors.
%%

\chapter{پروژه‌های بزرگ}
\label{chap:large-projects}

\section{چرا این موضوع مهم است؟}
\label{sec:large-projects-why}

وقتی کتاب شما به بیش از ۱۰۰ صفحه می‌رسد،
نوشتن همه‌ی محتوا در یک فایل \file{.tex}
غیرعملی می‌شود. کامپایل کند می‌شود، پیدا
کردن خطاها سخت می‌شود، و همکاری با دیگران
دشوار می‌شود.

\lr{MatinBook} به‌طور خاص برای پروژه‌های بزرگ
طراحی شده است. در این فصل، یاد می‌گیرید که
چگونه یک کتاب چندفایلی را سازمان‌دهی کنید.

\mnote{کتاب‌های حرفه‌ای معمولاً بیش از ۲۰۰ صفحه هستند و به‌طور طبیعی به چندین فایل تقسیم می‌شوند.}

\section{ساختار پیشنهادی}
\label{sec:large-projects-structure}

\subsection{ساختار ساده}
\label{sec:large-projects-simple}

برای کتاب‌های کوچک تا متوسط (کمتر از ۱۰
فصل):

\begin{latin}
\begin{minted}{text}
my-book/
├── main.tex
├── references.bib
├── chapters/
│   ├── chapter-01.tex
│   ├── chapter-02.tex
│   ├── chapter-03.tex
│   └── ...
├── figures/
│   ├── figure-01.pdf
│   └── ...
└── code/
    ├── example-01.py
    └── ...
\end{minted}
\end{latin}

\mnote{هر فصل در یک فایل جداگانه قرار می‌گیرد. این کار، کامپایل را سریع‌تر و مدیریت را آسان‌تر می‌کند.}

\subsection{ساختار پیشرفته}
\label{sec:large-projects-advanced}

برای کتاب‌های بزرگ (بیش از ۱۰ فصل یا
چند بخش):

\begin{latin}
\begin{minted}{text}
my-book/
├── main.tex
├── references.bib
├── frontmatter/
│   ├── cover.tex
│   ├── colophon.tex
│   ├── preface.tex
│   └── ...
├── part1-basics/
│   ├── chapter-01.tex
│   ├── chapter-02.tex
│   └── ...
├── part2-advanced/
│   ├── chapter-05.tex
│   ├── chapter-06.tex
│   └── ...
├── appendices/
│   ├── app-a.tex
│   └── ...
├── backmatter/
│   ├── bibliography.tex
│   └── index.tex
├── figures/
│   ├── chapter-01/
│   ├── chapter-02/
│   └── ...
└── code/
    ├── chapter-01/
    └── ...
\end{minted}
\end{latin}

\mnote{در ساختار پیشرفته، هر بخش از کتاب یک پوشه‌ی جداگانه دارد. این کار، مدیریت پروژه‌های بزرگ را آسان‌تر می‌کند.}

\section{فایل اصلی}
\label{sec:large-projects-main}

فایل \file{main.tex}، نقطه‌ی ورود پروژه
است. این فایل شامل:

\begin{itemize}
    \item \cmd{documentclass}.
    \item فراداده (\cmd{title}، \cmd{author}،
    \cmd{date}).
    \item \cmd{addbibresource}.
    \item \cmd{begin\{document\}}.
    \item جلد، پیش‌متن، متن اصلی، پس‌متن.
    \item \cmd{input} برای هر فصل.
    \item \cmd{end\{document\}}.
\end{itemize}

\begin{latin}
\begin{minted}{latex}
\documentclass{matinbook}

\addbibresource{references.bib}

\title{|\pc{عنوان کتاب}|}
\author{|\pc{نام نویسنده}|}
\date{|\pc{مهر ۱۴۰۵}|}

\booktitle{|\pc{عنوان کتاب}|}

\begin{document}

% |\pc{جلد جلو}|%
\makecover{...}{...}{...}{...}

% |\pc{پیش‌متن}|%
\frontmatter
\frontmattergeometry

\input{frontmatter/cover}
\input{frontmatter/preface}

% |\pc{متن اصلی}|%
\mainmatter
\mainmattergeometry

\part{|\pc{مبانی}|}
\input{part1-basics/chapter-01}
\input{part1-basics/chapter-02}

\part{|\pc{مباحث پیشرفته}|}
\input{part2-advanced/chapter-05}
\input{part2-advanced/chapter-06}

% |\pc{پس‌متن}|%
\backmatter
\frontmattergeometry

\printbibliography[title={|\pc{منابع و مراجع}|}]
\printindex

% |\pc{جلد عقب}|%
\makebackcover{...}

\end{document}
\end{minted}
\end{latin}

\mnote{فایل \lr{main.tex} باید کوتاه بماند. اگر بیش از ۱۰۰ خط شد، احتمالاً باید بخشی از محتوا را به فایل‌های جداگانه منتقل کنید.}

\section{فایل‌های فصل}
\label{sec:large-projects-chapters}

هر فایل فصل، فقط شامل محتوای همان فصل
است. \textbf{بدون} \cmd{documentclass} و
\textbf{بدون} \cmd{begin\{document\}}:

\begin{latin}
\begin{minted}{latex}
% |\pc{در فایل}|\lr{chapters/chapter-01.tex}:

\chapter{|\pc{مقدمه}|}
\label{chap:introduction}

\section{|\pc{انگیزه}|}
\label{sec:motivation}

|\pc{متن فصل...}|%

\begin{theorem}[|\pc{قضیه‌ی نمونه}|]
    |\pc{متن قضیه...}|%
    \label{thm:sample}
\end{theorem}

\begin{proof}
    |\pc{اثبات...}|%
\end{proof}
\end{minted}
\end{latin}

\mnote{فایل‌های فصل نباید \lr{\textbackslash documentclass} داشته باشند. آن‌ها فقط محتوا هستند و با \lr{\textbackslash input} وارد می‌شوند.}

\section{\cmd{input} در برابر \cmd{include}}
\label{sec:large-projects-input-include}

\lr{LaTeX} دو دستور برای وارد کردن فایل
ارائه می‌دهد:

\begin{itemize}
    \item \cmd{input\{file\}}: فایل را در
    همان‌جا وارد می‌کند.
    \item \cmd{include\{file\}}: فایل را
    در یک صفحه‌ی جدید وارد می‌کند.
\end{itemize}

\subsection{تفاوت‌ها}
\label{sec:large-projects-differences}

جدول \cref{tab:input-vs-include} تفاوت‌های
\cmd{input} و \cmd{include} را نشان می‌دهد.

\begin{matintable}{تفاوت \lr{input} و \lr{include}}{tab:input-vs-include}
\begin{tblr}{
    width = \textwidth,
    colspec = {l X[c] X[c]},
    hlines, vlines,
    colsep = 8pt,
    rowsep = 4pt,
    row{1} = {font=\bfseries},
}
ویژگی & \lr{\textbackslash input} & \lr{\textbackslash include} \\
صفحه‌ی جدید & خیر & بله \\
فایل \lr{.aux} جدا & خیر & بله \\
سرعت & سریع & کندتر \\
مناسب برای & فصل‌های کوتاه & فصل‌های بلند \\
\end{tblr}
\end{matintable}

\mnote{برای کتاب‌ها، معمولاً از \lr{\textbackslash input} استفاده می‌شود. \lr{\textbackslash include} فقط برای پروژه‌های بسیار بزرگ مفید است.}

\subsection{مثال}
\label{sec:large-projects-example}

\begin{latin}
\begin{minted}{latex}
% |\pc{با}|\lr{\textbackslash input}|%
\input{chapters/chapter-01}
\input{chapters/chapter-02}

% |\pc{با}|\lr{\textbackslash include}|%
\include{chapters/chapter-01}
\include{chapters/chapter-02}
\end{minted}
\end{latin}

\section{مدیریت شکل‌ها}
\label{sec:large-projects-figures}

برای پروژه‌های بزرگ، بهتر است شکل‌ها را در
پوشه‌های جداگانه سازمان‌دهی کنید:

\begin{latin}
\begin{minted}{text}
figures/
├── chapter-01/
│   ├── figure-01.pdf
│   ├── figure-02.pdf
│   └── ...
├── chapter-02/
│   └── ...
└── cover/
    └── cover-image.pdf
\end{minted}
\end{latin}

سپس در سند:

\begin{latin}
\begin{minted}{latex}
\includegraphics{figures/chapter-01/figure-01.pdf}
\end{minted}
\end{latin}

\mnote{برای کوتاه‌تر کردن مسیرها، می‌توانید از \lr{\textbackslash graphicspath} استفاده کنید که در \lr{MatinBook} از پیش تنظیم شده است.}

\section{مدیریت کد}
\label{sec:large-projects-code}

برای کدهای طولانی، بهتر است آن‌ها را در
فایل‌های جداگانه ذخیره کنید و با
\cmd{inputminted} وارد کنید:

\begin{latin}
\begin{minted}{text}
code/
├── chapter-01/
│   ├── example-01.py
│   └── example-02.py
└── chapter-02/
    └── ...
\end{minted}
\end{latin}

سپس در سند:

\begin{latin}
\begin{minted}{latex}
\begin{latin}
\inputminted{python}{code/chapter-01/example-01.py}
\end{latin}
\codecaption{|\pc{مثال اول}|}
\label{code:example-01}
\end{minted}
\end{latin}

\mnote{دستور \lr{\textbackslash inputminted} کد را از یک فایل خارجی وارد می‌کند. این کار، سورس کتاب را کوتاه‌تر می‌کند.}

\section{مدیریت کتاب‌نامه}
\label{sec:large-projects-bibliography}

برای کتاب‌های بزرگ، بهتر است کتاب‌نامه را
در چند فایل تقسیم کنید:

\begin{latin}
\begin{minted}{text}
bibliography/
├── books.bib
├── articles.bib
├── online.bib
└── persian.bib
\end{minted}
\end{latin}

سپس در \file{main.tex}:

\begin{latin}
\begin{minted}{latex}
\addbibresource{bibliography/books.bib}
\addbibresource{bibliography/articles.bib}
\addbibresource{bibliography/online.bib}
\addbibresource{bibliography/persian.bib}
\end{minted}
\end{latin}

\mnote{می‌توانید چند فایل \lr{.bib} را با \lr{\textbackslash addbibresource} بارگذاری کنید. \lr{biber} همه را با هم پردازش می‌کند.}

\section{همکاری با دیگران}
\label{sec:large-projects-collaboration}

اگر با دیگران روی یک کتاب کار می‌کنید،
این نکات را رعایت کنید:

\begin{enumerate}
    \item \textbf{استفاده از \lr{Git}:} نسخه‌بندی
    و همکاری.
    \item \textbf{تقسیم فایل‌ها:} هر نفر یک
    فصل.
    \item \textbf{قواعد نگارشی:} یکدست‌سازی
    سبک نوشتار.
    \item \textbf{برچسب‌های یکتا:} جلوگیری از
    تداخل برچسب‌ها.
    \item \textbf{کامپایل مکرر:} بررسی
    خطاها پس از هر تغییر.
\end{enumerate}

\mnote{در همکاری گروهی، \lr{Git} ابزار ضروری است. هر نفر روی برنچ خود کار می‌کند و سپس تغییرات را ادغام می‌کند.}

\section{کامپایل پروژه‌های بزرگ}
\label{sec:large-projects-compile}

\subsection{کامپایل کامل}
\label{sec:large-projects-compile-full}

\begin{latin}
\begin{minted}{bash}
# |\pc{مرحله‌ی اول: کامپایل اول}|%
xelatex -shell-escape main.tex

# |\pc{مرحله‌ی دوم: کتاب‌نامه}|%
biber main

# |\pc{مرحله‌ی سوم: کامپایل دوم}|%
xelatex -shell-escape main.tex

# |\pc{مرحله‌ی چهارم: کامپایل سوم}|%
xelatex -shell-escape main.tex
\end{minted}
\end{latin}

\mnote{برای پروژه‌های بزرگ، کامپایل کامل ممکن است چند دقیقه طول بکشد. از \lr{latexmk} برای خودکارسازی استفاده کنید.}

\subsection{استفاده از \lr{latexmk}}
\label{sec:large-projects-compile-latexmk}

\begin{latin}
\begin{minted}{bash}
latexmk -xelatex -shell-escape main.tex
\end{minted}
\end{latin}

\mnote{\lr{latexmk} به‌طور خودکار تعداد کامپایل‌های لازم را تشخیص می‌دهد و کل چرخه را اجرا می‌کند.}

\subsection{کامپایل فصل به فصل}
\label{sec:large-projects-compile-chapter}

برای صرفه‌جویی در زمان، می‌توانید فقط
یک فصل را کامپایل کنید:

\begin{latin}
\begin{minted}{bash}
# |\pc{در}|\lr{main.tex}|%
% |\pc{فقط فصل اول را کامپایل کن}|%
\includeonly{chapters/chapter-01}
\end{minted}
\end{latin}

\mnote{دستور \lr{\textbackslash includeonly} فقط با \lr{\textbackslash include} کار می‌کند، نه \lr{\textbackslash input}.}

\section{خطاهای رایج}
\label{sec:large-projects-mistakes}

\subsection{فراموش کردن \cmd{input}}
\label{sec:large-projects-mistake-input}

\textbf{مشکل:} اگر \cmd{input} را برای یک
فصل فراموش کنید، آن فصل در کتاب ظاهر
نمی‌شود.

\textbf{راه‌حل:} فهرست \cmd{input}ها را
با فهرست فصل‌ها مقایسه کنید.

\subsection{برچسب تکراری}
\label{sec:large-projects-mistake-duplicate}

\textbf{مشکل:} اگر در دو فصل برچسب یکسان
استفاده کنید، خطای \lr{Multiply defined labels}
می‌دهید.

\textbf{راه‌حل:} از پیشوندهای یکتا برای
هر فصل استفاده کنید (مثلاً \lr{ch01:}،
\lr{ch02:}).

\subsection{مسیر اشتباه}
\label{sec:large-projects-mistake-path}

\textbf{مشکل:} اگر مسیر فایل‌ها اشتباه
باشد، خطای \lr{File not found} می‌دهید.

\textbf{راه‌حل:} مسیرها را نسبت به
\file{main.tex} بررسی کنید.

\subsection{کامپایل ناقص}
\label{sec:large-projects-mistake-incomplete}

\textbf{مشکل:} اگر فقط یک بار کامپایل
کنید، ارجاعات و کتاب‌نامه کامل نمی‌شوند.

\textbf{راه‌حل:} چرخه‌ی کامل کامپایل را
اجرا کنید (حداقل ۲ بار \lr{xelatex} +
\lr{biber}).

\section{تمرین‌ها}
\label{sec:large-projects-exercises}

\begin{exercise}
    یک کتاب با ساختار پیشرفته (پوشه‌بندی
    شده) بسازید و آن را کامپایل کنید.
    \label{exc:large-projects-structure}
\end{exercise}

\begin{exercise}
    یک کتاب با دو بخش و پنج فصل بسازید.
    از \cmd{input} برای وارد کردن فصل‌ها
    استفاده کنید.
    \label{exc:large-projects-parts}
\end{exercise}

\begin{exercise}
    شکل‌ها را در پوشه‌های جداگانه سازمان‌دهی
    کنید و از \cmd{includegraphics} با
    مسیر کامل استفاده کنید.
    \label{exc:large-projects-figures}
\end{exercise}

\begin{exercise}
    یک کتاب با \cmd{includeonly} کامپایل
    کنید تا فقط یک فصل را ببینید.
    \label{exc:large-projects-includeonly}
\end{exercise}

\section{خلاصه}
\label{sec:large-projects-summary}

در این فصل، با مدیریت پروژه‌های بزرگ در
\lr{MatinBook} آشنا شدیم:

\begin{itemize}
    \item برای کتاب‌های بزرگ، محتوا را به
    فایل‌های جداگانه تقسیم کنید.

    \item ساختار ساده: \file{main.tex} +
    \file{chapters/} + \file{figures/} +
    \file{code/}.

    \item ساختار پیشرفته: تقسیم به بخش‌ها
    (\file{part1-*}، \file{part2-*} و...).

    \item \file{main.tex} باید کوتاه باشد
    (کمتر از ۱۰۰ خط).

    \item فایل‌های فصل نباید
    \cmd{documentclass} داشته باشند.

    \item \cmd{input} فایل را در همان‌جا
    وارد می‌کند؛ \cmd{include} صفحه‌ی
    جدید ایجاد می‌کند.

    \item برای کدهای طولانی، از
    \cmd{inputminted} استفاده کنید.

    \item کتاب‌نامه را می‌توان در چند فایل
    \lr{.bib} تقسیم کرد.

    \item در همکاری گروهی، از \lr{Git}
    و قواعد یکدست استفاده کنید.

    \item خطاهای رایج شامل فراموش کردن
    \cmd{input}، برچسب تکراری، مسیر
    اشتباه، و کامپایل ناقص است.
\end{itemize}

این فصل، آخرین فصل از بخش چهارم (مباحث
پیشرفته) بود. در بخش پنجم (مثال‌های کامل)،
سه فصل کامل از کتاب‌های ریاضی، برنامه‌نویسی
و فیزیک ارائه می‌شود.